% Red neuronal multifidelidad:
%   - red de baja fidelidad (BF) pre-entrenada, con pesos fijos, que usa solo el
%     subconjunto de entradas x_1..x_n;
%   - red de alta fidelidad (AF) informada por la fisica (PINN), que usa todas las
%     entradas x_1..x_N mas la prediccion u_BF.
% La columna de entradas se dibuja UNA sola vez y es compartida: las entradas del
% subconjunto (recuadro sombreado) salen dos veces, hacia la red BF y hacia el bus
% de todas las entradas, mientras que las restantes salen solo hacia el bus.
%
% Fragmento para \input directo (sin el paquete standalone).
% NO necesita ninguna libreria de TikZ: solo \usepackage{tikz} y \usepackage{amsmath}.
% Uso tipico:
%   \begin{figure}[htbp]\centering
%     \resizebox{\textwidth}{!}{% Red neuronal multifidelidad:
%   - red de baja fidelidad (BF) pre-entrenada, con pesos fijos, que usa solo el
%     subconjunto de entradas x_1..x_n;
%   - red de alta fidelidad (AF) informada por la fisica (PINN), que usa todas las
%     entradas x_1..x_N mas la prediccion u_BF.
% La columna de entradas se dibuja UNA sola vez y es compartida: las entradas del
% subconjunto (recuadro sombreado) salen dos veces, hacia la red BF y hacia el bus
% de todas las entradas, mientras que las restantes salen solo hacia el bus.
%
% Fragmento para \input directo (sin el paquete standalone).
% NO necesita ninguna libreria de TikZ: solo \usepackage{tikz} y \usepackage{amsmath}.
% Uso tipico:
%   \begin{figure}[htbp]\centering
%     \resizebox{\textwidth}{!}{% Red neuronal multifidelidad:
%   - red de baja fidelidad (BF) pre-entrenada, con pesos fijos, que usa solo el
%     subconjunto de entradas x_1..x_n;
%   - red de alta fidelidad (AF) informada por la fisica (PINN), que usa todas las
%     entradas x_1..x_N mas la prediccion u_BF.
% La columna de entradas se dibuja UNA sola vez y es compartida: las entradas del
% subconjunto (recuadro sombreado) salen dos veces, hacia la red BF y hacia el bus
% de todas las entradas, mientras que las restantes salen solo hacia el bus.
%
% Fragmento para \input directo (sin el paquete standalone).
% NO necesita ninguna libreria de TikZ: solo \usepackage{tikz} y \usepackage{amsmath}.
% Uso tipico:
%   \begin{figure}[htbp]\centering
%     \resizebox{\textwidth}{!}{% Red neuronal multifidelidad:
%   - red de baja fidelidad (BF) pre-entrenada, con pesos fijos, que usa solo el
%     subconjunto de entradas x_1..x_n;
%   - red de alta fidelidad (AF) informada por la fisica (PINN), que usa todas las
%     entradas x_1..x_N mas la prediccion u_BF.
% La columna de entradas se dibuja UNA sola vez y es compartida: las entradas del
% subconjunto (recuadro sombreado) salen dos veces, hacia la red BF y hacia el bus
% de todas las entradas, mientras que las restantes salen solo hacia el bus.
%
% Fragmento para \input directo (sin el paquete standalone).
% NO necesita ninguna libreria de TikZ: solo \usepackage{tikz} y \usepackage{amsmath}.
% Uso tipico:
%   \begin{figure}[htbp]\centering
%     \resizebox{\textwidth}{!}{\input{Figures/mfnn_pinn.tex}}
%     \caption{...}\label{fig:mfnn_pinn}
%   \end{figure}
%
% Notas de implementacion (todas para que el fragmento no dependa del preambulo):
%  - Sin \usetikzlibrary aqui adentro: ejecutado en modo horizontal (dentro de un
%    \resizebox) deja ~55pt de espacios que ensanchan la caja del \input y hacen que
%    \resizebox empequenezca el dibujo; y no se puede encerrar en un grupo o \hbox
%    para absorberlos porque las definiciones de las librerias son locales.
%  - Sin la libreria backgrounds: los recuadros de las entradas se dibujan ANTES que
%    los circulos, de modo que el relleno queda debajo sin necesidad de capas.
%  - Sin la libreria fit: los recuadros de las redes se dan por coordenadas explicitas.
%  - Sin la libreria arrows.meta: las puntas son `stealth`, que ya trae pgf. Ademas se
%    escriben como -stealth y no con la clave >=, para no depender de que babel-spanish
%    tenga desactivado el shorthand `>`.
%  - Los tamanos de fuente son absolutos (\fontsize) para que la geometria no dependa
%    del tamano base del documento que hace el \input.
%  - El dibujo es casi cuadrado (15.3 x 15.3 cm) para que al escalarlo al ancho del
%    texto las etiquetas no queden diminutas.
\begin{figure}[hbtp]
  \centering
  \resizebox{\textwidth}{!}{
\begin{tikzpicture}[
  chico/.style={font=\fontsize{9}{11}\selectfont},
  % un unico estilo de circulo: entradas, neuronas y salidas miden todas lo mismo
  entrada/.style={circle,draw=black!70,fill=white,minimum size=0.9cm,inner sep=0pt,
                  font=\fontsize{9}{11}\selectfont},
  entradaG/.style={entrada},
  neurona/.style={entrada},
  salida/.style={entrada},
  syn/.style={draw=black!25,line width=0.3pt},
  reparto/.style={draw=black!45,line width=0.4pt},
  flujo/.style={-stealth,draw=black!75,line width=0.8pt},
  tronco/.style={draw=black!75,line width=0.8pt},
  bus/.style={-stealth,draw=black!75,line width=1.5pt},
  bloque/.style={rounded corners=2pt,draw=black!70,align=center,inner sep=5pt,
                 font=\fontsize{9}{11}\selectfont,text width=3.2cm},
  perdida/.style={bloque,fill=black!6},
  barra/.style={rounded corners=2pt,draw=black!70,fill=black!10,
                minimum width=0.6cm,minimum height=4.4cm,inner sep=0pt},
  grupoSub/.style={rounded corners=4pt,draw=black!40,fill=black!7},
  grupoResto/.style={rounded corners=4pt,draw=black!40,dash pattern=on 2pt off 1.5pt},
  grupoBF/.style={draw=black!45,dashed,rounded corners=6pt},
  grupoAF/.style={draw=black!70,rounded corners=6pt},
  etiqueta/.style={font=\fontsize{9}{11}\bfseries\selectfont,align=left,text width=5.2cm},
  nota/.style={font=\fontsize{8}{9.6}\selectfont,align=left},
  notac/.style={font=\fontsize{8}{9.6}\selectfont,align=center},
  bp/.style={-stealth,dashed,draw=black!60,line width=0.8pt},
]

%=============== COLUMNA DE ENTRADAS (UNICA Y COMPARTIDA) ==============
% los recuadros van primero para que el relleno quede por debajo de los circulos
\draw[grupoSub]   (-0.66, 0.09) rectangle (0.66, 4.06);
\draw[grupoResto] (-0.66,-3.21) rectangle (0.66,-0.39);

% subconjunto que ve la red de baja fidelidad
\node[entrada]  (e1)  at (0, 3.40) {$x_1$};
\node[entrada]  (e2)  at (0, 2.35) {$x_2$};
\node[chico]    (ev1) at (0, 1.55) {$\vdots$};
\node[entrada]  (en)  at (0, 0.75) {$x_n$};
% entradas restantes, exclusivas de la red de alta fidelidad
\node[entradaG] (er1) at (0,-1.05) {$x_{n+1}$};
\node[chico]    (ev2) at (0,-1.80) {$\vdots$};
\node[entrada]  (eN)  at (0,-2.55) {$x_N$};

% \node[notac,rotate=90] at (-1.35, 2.02) {primeras $n$\\ entradas};
% \node[notac,rotate=90] at (-1.35,-1.80) {entradas\\ restantes};

%=================== RED DE BAJA FIDELIDAD (BF) ========================
\foreach \j/\y in {1/3.07,2/2.02,3/0.97} {
  \node[neurona] (bf1-\j) at (2.6,\y) {};
  \node[neurona] (bf2-\j) at (4.1,\y) {};
}
\node[salida] (ubf) at (5.6,2.02) {$u_\text{LF}$};

% solo el subconjunto alimenta la red BF
\foreach \e in {e1,e2,en} {
  \foreach \j in {1,2,3} { \draw[syn] (\e) -- (bf1-\j); }
}
\foreach \i in {1,2,3} {
  \foreach \j in {1,2,3} { \draw[syn] (bf1-\i) -- (bf2-\j); }
  \draw[syn] (bf2-\i) -- (ubf);
}

\draw[grupoBF] (1.83,0.20) rectangle (6.37,3.84);
\node[etiqueta,anchor=south west] at (1.83,3.84)
  {Red de baja fidelidad $\theta_\text{LF}$ pre-entrenada};

%=========== BUS CON TODAS LAS ENTRADAS HACIA LA RED AF ================
% cada entrada, del subconjunto y de las restantes, se reparte al bus
\coordinate (nodoBus) at (2.6,-3.0);
\foreach \e in {e1,e2,en,er1,eN} { \draw[reparto] (\e) -- (nodoBus); }
\fill[black!75] (nodoBus) circle (1.6pt);

\draw[bus,rounded corners=3pt] (nodoBus) -- (5.9,-3.0) -- (5.9,-1.2) -- (7.30,-1.2);
\node[nota,anchor=north west] at (2.9,-3.4) {todas las entradas\\ $x_1,\dots,x_N$};

%=================== RED DE ALTA FIDELIDAD (AF) ========================
\node[barra] (concat) at (7.6,0) {};
\node[nota,rotate=90] at (7.6,0) {concatenación};

\foreach \j/\y in {1/1.70,2/0.57,3/-0.57,4/-1.70} {
  \node[neurona] (af1-\j) at (8.9,\y) {};
  \node[neurona] (af2-\j) at (10.4,\y) {};
}
\node[salida] (uaf) at (11.9,0) {$u_\text{HF}$};

\draw[flujo] (ubf.east) -- (7.30,1.3);
\foreach \j in {1,2,3,4} { \draw[syn] (concat.east) -- (af1-\j); }
\foreach \i in {1,2,3,4} {
  \foreach \j in {1,2,3,4} { \draw[syn] (af1-\i) -- (af2-\j); }
  \draw[syn] (af2-\i) -- (uaf);
}

\draw[grupoAF] (6.98,-2.52) rectangle (12.67,2.52);
\node[etiqueta,anchor=south west] at (6.98,2.52)
  {Red de alta fidelidad $\theta_\text{HF}$\\ informada por la física (PINN)};

%=============== RAMAS DE PÉRDIDA (FÍSICA Y DATOS) =====================
\node[bloque]  (ad)   at (4.6,-6.3)
  {Derivadas\\ automáticas\\ de $u_\text{HF}$ respecto\\ de las entradas};
\node[perdida] (lfis) at (4.6,-8.4)
  {$\mathcal{L}_{\text{fís}}=\big\|\mathcal{R}(u_\text{HF})\big\|^{2}$\\ residuo de la EDP};
\node[perdida,text width=3.4cm] (ldat) at (8.8,-6.3)
  {$\mathcal{L}_{\text{datos}}=\big\|u_\text{HF}-u^{*}\big\|^{2}$\\ $u^{*}$: datos medidos};
\node[perdida,text width=3.4cm] (ltot) at (8.8,-8.4)
  {$\mathcal{L}=\mathcal{L}_{\text{datos}}+\lambda\,\mathcal{L}_{\text{fís}}$};

% tronco desde u_AF, con derivacion hacia la rama de datos
\coordinate (bif) at (11.9,-4.9);
\draw[tronco] (uaf.south) -- (bif);
\draw[flujo,rounded corners=3pt] (bif) -- (4.6,-4.9) -- (ad.north);
\fill[black!75] (8.8,-4.9) circle (1.6pt);
\draw[flujo] (8.8,-4.9) -- (ldat.north);

\draw[flujo] (ad.south)   -- (lfis.north);
\draw[flujo] (ldat.south) -- (ltot.north);
\draw[flujo] (lfis.east)  -- (ltot.west);

%=================== RETROPROPAGACIÓN (SOLO RED AF) ====================
% vuelve por la derecha y entra al costado de la red AF, sin cruzar el tronco
\draw[bp,rounded corners=3pt]
  (ltot.east) -- (13.4,-8.4) -- (13.4,-1.8) -- (12.67,-1.8);
% \node[notac,anchor=north,text width=4.2cm] at (8.8,-9.3)
%   {retropropagación: sólo se\\ actualizan los pesos $\theta_{AF}$};

%============================ ACLARACIÓN ===============================
% \node[nota,anchor=north west,text width=4.0cm] at (-1.75,-4.9)
%   {la red BF sólo usa $x_1,\dots,x_n$; la red AF usa todas las entradas y
%    además $u_{BF}$};

\end{tikzpicture}
  }
  \caption{Esquema de la arquitectura de red neuronal de multifidelidad informada por física.}
  \label{fig:mfpinn}
\end{figure}
}
%     \caption{...}\label{fig:mfnn_pinn}
%   \end{figure}
%
% Notas de implementacion (todas para que el fragmento no dependa del preambulo):
%  - Sin \usetikzlibrary aqui adentro: ejecutado en modo horizontal (dentro de un
%    \resizebox) deja ~55pt de espacios que ensanchan la caja del \input y hacen que
%    \resizebox empequenezca el dibujo; y no se puede encerrar en un grupo o \hbox
%    para absorberlos porque las definiciones de las librerias son locales.
%  - Sin la libreria backgrounds: los recuadros de las entradas se dibujan ANTES que
%    los circulos, de modo que el relleno queda debajo sin necesidad de capas.
%  - Sin la libreria fit: los recuadros de las redes se dan por coordenadas explicitas.
%  - Sin la libreria arrows.meta: las puntas son `stealth`, que ya trae pgf. Ademas se
%    escriben como -stealth y no con la clave >=, para no depender de que babel-spanish
%    tenga desactivado el shorthand `>`.
%  - Los tamanos de fuente son absolutos (\fontsize) para que la geometria no dependa
%    del tamano base del documento que hace el \input.
%  - El dibujo es casi cuadrado (15.3 x 15.3 cm) para que al escalarlo al ancho del
%    texto las etiquetas no queden diminutas.
\begin{figure}[hbtp]
  \centering
  \resizebox{\textwidth}{!}{
\begin{tikzpicture}[
  chico/.style={font=\fontsize{9}{11}\selectfont},
  % un unico estilo de circulo: entradas, neuronas y salidas miden todas lo mismo
  entrada/.style={circle,draw=black!70,fill=white,minimum size=0.9cm,inner sep=0pt,
                  font=\fontsize{9}{11}\selectfont},
  entradaG/.style={entrada},
  neurona/.style={entrada},
  salida/.style={entrada},
  syn/.style={draw=black!25,line width=0.3pt},
  reparto/.style={draw=black!45,line width=0.4pt},
  flujo/.style={-stealth,draw=black!75,line width=0.8pt},
  tronco/.style={draw=black!75,line width=0.8pt},
  bus/.style={-stealth,draw=black!75,line width=1.5pt},
  bloque/.style={rounded corners=2pt,draw=black!70,align=center,inner sep=5pt,
                 font=\fontsize{9}{11}\selectfont,text width=3.2cm},
  perdida/.style={bloque,fill=black!6},
  barra/.style={rounded corners=2pt,draw=black!70,fill=black!10,
                minimum width=0.6cm,minimum height=4.4cm,inner sep=0pt},
  grupoSub/.style={rounded corners=4pt,draw=black!40,fill=black!7},
  grupoResto/.style={rounded corners=4pt,draw=black!40,dash pattern=on 2pt off 1.5pt},
  grupoBF/.style={draw=black!45,dashed,rounded corners=6pt},
  grupoAF/.style={draw=black!70,rounded corners=6pt},
  etiqueta/.style={font=\fontsize{9}{11}\bfseries\selectfont,align=left,text width=5.2cm},
  nota/.style={font=\fontsize{8}{9.6}\selectfont,align=left},
  notac/.style={font=\fontsize{8}{9.6}\selectfont,align=center},
  bp/.style={-stealth,dashed,draw=black!60,line width=0.8pt},
]

%=============== COLUMNA DE ENTRADAS (UNICA Y COMPARTIDA) ==============
% los recuadros van primero para que el relleno quede por debajo de los circulos
\draw[grupoSub]   (-0.66, 0.09) rectangle (0.66, 4.06);
\draw[grupoResto] (-0.66,-3.21) rectangle (0.66,-0.39);

% subconjunto que ve la red de baja fidelidad
\node[entrada]  (e1)  at (0, 3.40) {$x_1$};
\node[entrada]  (e2)  at (0, 2.35) {$x_2$};
\node[chico]    (ev1) at (0, 1.55) {$\vdots$};
\node[entrada]  (en)  at (0, 0.75) {$x_n$};
% entradas restantes, exclusivas de la red de alta fidelidad
\node[entradaG] (er1) at (0,-1.05) {$x_{n+1}$};
\node[chico]    (ev2) at (0,-1.80) {$\vdots$};
\node[entrada]  (eN)  at (0,-2.55) {$x_N$};

% \node[notac,rotate=90] at (-1.35, 2.02) {primeras $n$\\ entradas};
% \node[notac,rotate=90] at (-1.35,-1.80) {entradas\\ restantes};

%=================== RED DE BAJA FIDELIDAD (BF) ========================
\foreach \j/\y in {1/3.07,2/2.02,3/0.97} {
  \node[neurona] (bf1-\j) at (2.6,\y) {};
  \node[neurona] (bf2-\j) at (4.1,\y) {};
}
\node[salida] (ubf) at (5.6,2.02) {$u_\text{LF}$};

% solo el subconjunto alimenta la red BF
\foreach \e in {e1,e2,en} {
  \foreach \j in {1,2,3} { \draw[syn] (\e) -- (bf1-\j); }
}
\foreach \i in {1,2,3} {
  \foreach \j in {1,2,3} { \draw[syn] (bf1-\i) -- (bf2-\j); }
  \draw[syn] (bf2-\i) -- (ubf);
}

\draw[grupoBF] (1.83,0.20) rectangle (6.37,3.84);
\node[etiqueta,anchor=south west] at (1.83,3.84)
  {Red de baja fidelidad $\theta_\text{LF}$ pre-entrenada};

%=========== BUS CON TODAS LAS ENTRADAS HACIA LA RED AF ================
% cada entrada, del subconjunto y de las restantes, se reparte al bus
\coordinate (nodoBus) at (2.6,-3.0);
\foreach \e in {e1,e2,en,er1,eN} { \draw[reparto] (\e) -- (nodoBus); }
\fill[black!75] (nodoBus) circle (1.6pt);

\draw[bus,rounded corners=3pt] (nodoBus) -- (5.9,-3.0) -- (5.9,-1.2) -- (7.30,-1.2);
\node[nota,anchor=north west] at (2.9,-3.4) {todas las entradas\\ $x_1,\dots,x_N$};

%=================== RED DE ALTA FIDELIDAD (AF) ========================
\node[barra] (concat) at (7.6,0) {};
\node[nota,rotate=90] at (7.6,0) {concatenación};

\foreach \j/\y in {1/1.70,2/0.57,3/-0.57,4/-1.70} {
  \node[neurona] (af1-\j) at (8.9,\y) {};
  \node[neurona] (af2-\j) at (10.4,\y) {};
}
\node[salida] (uaf) at (11.9,0) {$u_\text{HF}$};

\draw[flujo] (ubf.east) -- (7.30,1.3);
\foreach \j in {1,2,3,4} { \draw[syn] (concat.east) -- (af1-\j); }
\foreach \i in {1,2,3,4} {
  \foreach \j in {1,2,3,4} { \draw[syn] (af1-\i) -- (af2-\j); }
  \draw[syn] (af2-\i) -- (uaf);
}

\draw[grupoAF] (6.98,-2.52) rectangle (12.67,2.52);
\node[etiqueta,anchor=south west] at (6.98,2.52)
  {Red de alta fidelidad $\theta_\text{HF}$\\ informada por la física (PINN)};

%=============== RAMAS DE PÉRDIDA (FÍSICA Y DATOS) =====================
\node[bloque]  (ad)   at (4.6,-6.3)
  {Derivadas\\ automáticas\\ de $u_\text{HF}$ respecto\\ de las entradas};
\node[perdida] (lfis) at (4.6,-8.4)
  {$\mathcal{L}_{\text{fís}}=\big\|\mathcal{R}(u_\text{HF})\big\|^{2}$\\ residuo de la EDP};
\node[perdida,text width=3.4cm] (ldat) at (8.8,-6.3)
  {$\mathcal{L}_{\text{datos}}=\big\|u_\text{HF}-u^{*}\big\|^{2}$\\ $u^{*}$: datos medidos};
\node[perdida,text width=3.4cm] (ltot) at (8.8,-8.4)
  {$\mathcal{L}=\mathcal{L}_{\text{datos}}+\lambda\,\mathcal{L}_{\text{fís}}$};

% tronco desde u_AF, con derivacion hacia la rama de datos
\coordinate (bif) at (11.9,-4.9);
\draw[tronco] (uaf.south) -- (bif);
\draw[flujo,rounded corners=3pt] (bif) -- (4.6,-4.9) -- (ad.north);
\fill[black!75] (8.8,-4.9) circle (1.6pt);
\draw[flujo] (8.8,-4.9) -- (ldat.north);

\draw[flujo] (ad.south)   -- (lfis.north);
\draw[flujo] (ldat.south) -- (ltot.north);
\draw[flujo] (lfis.east)  -- (ltot.west);

%=================== RETROPROPAGACIÓN (SOLO RED AF) ====================
% vuelve por la derecha y entra al costado de la red AF, sin cruzar el tronco
\draw[bp,rounded corners=3pt]
  (ltot.east) -- (13.4,-8.4) -- (13.4,-1.8) -- (12.67,-1.8);
% \node[notac,anchor=north,text width=4.2cm] at (8.8,-9.3)
%   {retropropagación: sólo se\\ actualizan los pesos $\theta_{AF}$};

%============================ ACLARACIÓN ===============================
% \node[nota,anchor=north west,text width=4.0cm] at (-1.75,-4.9)
%   {la red BF sólo usa $x_1,\dots,x_n$; la red AF usa todas las entradas y
%    además $u_{BF}$};

\end{tikzpicture}
  }
  \caption{Esquema de la arquitectura de red neuronal de multifidelidad informada por física.}
  \label{fig:mfpinn}
\end{figure}
}
%     \caption{...}\label{fig:mfnn_pinn}
%   \end{figure}
%
% Notas de implementacion (todas para que el fragmento no dependa del preambulo):
%  - Sin \usetikzlibrary aqui adentro: ejecutado en modo horizontal (dentro de un
%    \resizebox) deja ~55pt de espacios que ensanchan la caja del \input y hacen que
%    \resizebox empequenezca el dibujo; y no se puede encerrar en un grupo o \hbox
%    para absorberlos porque las definiciones de las librerias son locales.
%  - Sin la libreria backgrounds: los recuadros de las entradas se dibujan ANTES que
%    los circulos, de modo que el relleno queda debajo sin necesidad de capas.
%  - Sin la libreria fit: los recuadros de las redes se dan por coordenadas explicitas.
%  - Sin la libreria arrows.meta: las puntas son `stealth`, que ya trae pgf. Ademas se
%    escriben como -stealth y no con la clave >=, para no depender de que babel-spanish
%    tenga desactivado el shorthand `>`.
%  - Los tamanos de fuente son absolutos (\fontsize) para que la geometria no dependa
%    del tamano base del documento que hace el \input.
%  - El dibujo es casi cuadrado (15.3 x 15.3 cm) para que al escalarlo al ancho del
%    texto las etiquetas no queden diminutas.
\begin{figure}[hbtp]
  \centering
  \resizebox{\textwidth}{!}{
\begin{tikzpicture}[
  chico/.style={font=\fontsize{9}{11}\selectfont},
  % un unico estilo de circulo: entradas, neuronas y salidas miden todas lo mismo
  entrada/.style={circle,draw=black!70,fill=white,minimum size=0.9cm,inner sep=0pt,
                  font=\fontsize{9}{11}\selectfont},
  entradaG/.style={entrada},
  neurona/.style={entrada},
  salida/.style={entrada},
  syn/.style={draw=black!25,line width=0.3pt},
  reparto/.style={draw=black!45,line width=0.4pt},
  flujo/.style={-stealth,draw=black!75,line width=0.8pt},
  tronco/.style={draw=black!75,line width=0.8pt},
  bus/.style={-stealth,draw=black!75,line width=1.5pt},
  bloque/.style={rounded corners=2pt,draw=black!70,align=center,inner sep=5pt,
                 font=\fontsize{9}{11}\selectfont,text width=3.2cm},
  perdida/.style={bloque,fill=black!6},
  barra/.style={rounded corners=2pt,draw=black!70,fill=black!10,
                minimum width=0.6cm,minimum height=4.4cm,inner sep=0pt},
  grupoSub/.style={rounded corners=4pt,draw=black!40,fill=black!7},
  grupoResto/.style={rounded corners=4pt,draw=black!40,dash pattern=on 2pt off 1.5pt},
  grupoBF/.style={draw=black!45,dashed,rounded corners=6pt},
  grupoAF/.style={draw=black!70,rounded corners=6pt},
  etiqueta/.style={font=\fontsize{9}{11}\bfseries\selectfont,align=left,text width=5.2cm},
  nota/.style={font=\fontsize{8}{9.6}\selectfont,align=left},
  notac/.style={font=\fontsize{8}{9.6}\selectfont,align=center},
  bp/.style={-stealth,dashed,draw=black!60,line width=0.8pt},
]

%=============== COLUMNA DE ENTRADAS (UNICA Y COMPARTIDA) ==============
% los recuadros van primero para que el relleno quede por debajo de los circulos
\draw[grupoSub]   (-0.66, 0.09) rectangle (0.66, 4.06);
\draw[grupoResto] (-0.66,-3.21) rectangle (0.66,-0.39);

% subconjunto que ve la red de baja fidelidad
\node[entrada]  (e1)  at (0, 3.40) {$x_1$};
\node[entrada]  (e2)  at (0, 2.35) {$x_2$};
\node[chico]    (ev1) at (0, 1.55) {$\vdots$};
\node[entrada]  (en)  at (0, 0.75) {$x_n$};
% entradas restantes, exclusivas de la red de alta fidelidad
\node[entradaG] (er1) at (0,-1.05) {$x_{n+1}$};
\node[chico]    (ev2) at (0,-1.80) {$\vdots$};
\node[entrada]  (eN)  at (0,-2.55) {$x_N$};

% \node[notac,rotate=90] at (-1.35, 2.02) {primeras $n$\\ entradas};
% \node[notac,rotate=90] at (-1.35,-1.80) {entradas\\ restantes};

%=================== RED DE BAJA FIDELIDAD (BF) ========================
\foreach \j/\y in {1/3.07,2/2.02,3/0.97} {
  \node[neurona] (bf1-\j) at (2.6,\y) {};
  \node[neurona] (bf2-\j) at (4.1,\y) {};
}
\node[salida] (ubf) at (5.6,2.02) {$u_\text{LF}$};

% solo el subconjunto alimenta la red BF
\foreach \e in {e1,e2,en} {
  \foreach \j in {1,2,3} { \draw[syn] (\e) -- (bf1-\j); }
}
\foreach \i in {1,2,3} {
  \foreach \j in {1,2,3} { \draw[syn] (bf1-\i) -- (bf2-\j); }
  \draw[syn] (bf2-\i) -- (ubf);
}

\draw[grupoBF] (1.83,0.20) rectangle (6.37,3.84);
\node[etiqueta,anchor=south west] at (1.83,3.84)
  {Red de baja fidelidad $\theta_\text{LF}$ pre-entrenada};

%=========== BUS CON TODAS LAS ENTRADAS HACIA LA RED AF ================
% cada entrada, del subconjunto y de las restantes, se reparte al bus
\coordinate (nodoBus) at (2.6,-3.0);
\foreach \e in {e1,e2,en,er1,eN} { \draw[reparto] (\e) -- (nodoBus); }
\fill[black!75] (nodoBus) circle (1.6pt);

\draw[bus,rounded corners=3pt] (nodoBus) -- (5.9,-3.0) -- (5.9,-1.2) -- (7.30,-1.2);
\node[nota,anchor=north west] at (2.9,-3.4) {todas las entradas\\ $x_1,\dots,x_N$};

%=================== RED DE ALTA FIDELIDAD (AF) ========================
\node[barra] (concat) at (7.6,0) {};
\node[nota,rotate=90] at (7.6,0) {concatenación};

\foreach \j/\y in {1/1.70,2/0.57,3/-0.57,4/-1.70} {
  \node[neurona] (af1-\j) at (8.9,\y) {};
  \node[neurona] (af2-\j) at (10.4,\y) {};
}
\node[salida] (uaf) at (11.9,0) {$u_\text{HF}$};

\draw[flujo] (ubf.east) -- (7.30,1.3);
\foreach \j in {1,2,3,4} { \draw[syn] (concat.east) -- (af1-\j); }
\foreach \i in {1,2,3,4} {
  \foreach \j in {1,2,3,4} { \draw[syn] (af1-\i) -- (af2-\j); }
  \draw[syn] (af2-\i) -- (uaf);
}

\draw[grupoAF] (6.98,-2.52) rectangle (12.67,2.52);
\node[etiqueta,anchor=south west] at (6.98,2.52)
  {Red de alta fidelidad $\theta_\text{HF}$\\ informada por la física (PINN)};

%=============== RAMAS DE PÉRDIDA (FÍSICA Y DATOS) =====================
\node[bloque]  (ad)   at (4.6,-6.3)
  {Derivadas\\ automáticas\\ de $u_\text{HF}$ respecto\\ de las entradas};
\node[perdida] (lfis) at (4.6,-8.4)
  {$\mathcal{L}_{\text{fís}}=\big\|\mathcal{R}(u_\text{HF})\big\|^{2}$\\ residuo de la EDP};
\node[perdida,text width=3.4cm] (ldat) at (8.8,-6.3)
  {$\mathcal{L}_{\text{datos}}=\big\|u_\text{HF}-u^{*}\big\|^{2}$\\ $u^{*}$: datos medidos};
\node[perdida,text width=3.4cm] (ltot) at (8.8,-8.4)
  {$\mathcal{L}=\mathcal{L}_{\text{datos}}+\lambda\,\mathcal{L}_{\text{fís}}$};

% tronco desde u_AF, con derivacion hacia la rama de datos
\coordinate (bif) at (11.9,-4.9);
\draw[tronco] (uaf.south) -- (bif);
\draw[flujo,rounded corners=3pt] (bif) -- (4.6,-4.9) -- (ad.north);
\fill[black!75] (8.8,-4.9) circle (1.6pt);
\draw[flujo] (8.8,-4.9) -- (ldat.north);

\draw[flujo] (ad.south)   -- (lfis.north);
\draw[flujo] (ldat.south) -- (ltot.north);
\draw[flujo] (lfis.east)  -- (ltot.west);

%=================== RETROPROPAGACIÓN (SOLO RED AF) ====================
% vuelve por la derecha y entra al costado de la red AF, sin cruzar el tronco
\draw[bp,rounded corners=3pt]
  (ltot.east) -- (13.4,-8.4) -- (13.4,-1.8) -- (12.67,-1.8);
% \node[notac,anchor=north,text width=4.2cm] at (8.8,-9.3)
%   {retropropagación: sólo se\\ actualizan los pesos $\theta_{AF}$};

%============================ ACLARACIÓN ===============================
% \node[nota,anchor=north west,text width=4.0cm] at (-1.75,-4.9)
%   {la red BF sólo usa $x_1,\dots,x_n$; la red AF usa todas las entradas y
%    además $u_{BF}$};

\end{tikzpicture}
  }
  \caption{Esquema de la arquitectura de red neuronal de multifidelidad informada por física.}
  \label{fig:mfpinn}
\end{figure}
}
%     \caption{...}\label{fig:mfnn_pinn}
%   \end{figure}
%
% Notas de implementacion (todas para que el fragmento no dependa del preambulo):
%  - Sin \usetikzlibrary aqui adentro: ejecutado en modo horizontal (dentro de un
%    \resizebox) deja ~55pt de espacios que ensanchan la caja del \input y hacen que
%    \resizebox empequenezca el dibujo; y no se puede encerrar en un grupo o \hbox
%    para absorberlos porque las definiciones de las librerias son locales.
%  - Sin la libreria backgrounds: los recuadros de las entradas se dibujan ANTES que
%    los circulos, de modo que el relleno queda debajo sin necesidad de capas.
%  - Sin la libreria fit: los recuadros de las redes se dan por coordenadas explicitas.
%  - Sin la libreria arrows.meta: las puntas son `stealth`, que ya trae pgf. Ademas se
%    escriben como -stealth y no con la clave >=, para no depender de que babel-spanish
%    tenga desactivado el shorthand `>`.
%  - Los tamanos de fuente son absolutos (\fontsize) para que la geometria no dependa
%    del tamano base del documento que hace el \input.
%  - El dibujo es casi cuadrado (15.3 x 15.3 cm) para que al escalarlo al ancho del
%    texto las etiquetas no queden diminutas.
\begin{figure}[hbtp]
  \centering
  \resizebox{\textwidth}{!}{
\begin{tikzpicture}[
  chico/.style={font=\fontsize{9}{11}\selectfont},
  % un unico estilo de circulo: entradas, neuronas y salidas miden todas lo mismo
  entrada/.style={circle,draw=black!70,fill=white,minimum size=0.9cm,inner sep=0pt,
                  font=\fontsize{9}{11}\selectfont},
  entradaG/.style={entrada},
  neurona/.style={entrada},
  salida/.style={entrada},
  syn/.style={draw=black!25,line width=0.3pt},
  reparto/.style={draw=black!45,line width=0.4pt},
  flujo/.style={-stealth,draw=black!75,line width=0.8pt},
  tronco/.style={draw=black!75,line width=0.8pt},
  bus/.style={-stealth,draw=black!75,line width=1.5pt},
  bloque/.style={rounded corners=2pt,draw=black!70,align=center,inner sep=5pt,
                 font=\fontsize{9}{11}\selectfont,text width=3.2cm},
  perdida/.style={bloque,fill=black!6},
  barra/.style={rounded corners=2pt,draw=black!70,fill=black!10,
                minimum width=0.6cm,minimum height=4.4cm,inner sep=0pt},
  grupoSub/.style={rounded corners=4pt,draw=black!40,fill=black!7},
  grupoResto/.style={rounded corners=4pt,draw=black!40,dash pattern=on 2pt off 1.5pt},
  grupoBF/.style={draw=black!45,dashed,rounded corners=6pt},
  grupoAF/.style={draw=black!70,rounded corners=6pt},
  etiqueta/.style={font=\fontsize{9}{11}\bfseries\selectfont,align=left,text width=5.2cm},
  nota/.style={font=\fontsize{8}{9.6}\selectfont,align=left},
  notac/.style={font=\fontsize{8}{9.6}\selectfont,align=center},
  bp/.style={-stealth,dashed,draw=black!60,line width=0.8pt},
]

%=============== COLUMNA DE ENTRADAS (UNICA Y COMPARTIDA) ==============
% los recuadros van primero para que el relleno quede por debajo de los circulos
\draw[grupoSub]   (-0.66, 0.09) rectangle (0.66, 4.06);
\draw[grupoResto] (-0.66,-3.21) rectangle (0.66,-0.39);

% subconjunto que ve la red de baja fidelidad
\node[entrada]  (e1)  at (0, 3.40) {$x_1$};
\node[entrada]  (e2)  at (0, 2.35) {$x_2$};
\node[chico]    (ev1) at (0, 1.55) {$\vdots$};
\node[entrada]  (en)  at (0, 0.75) {$x_n$};
% entradas restantes, exclusivas de la red de alta fidelidad
\node[entradaG] (er1) at (0,-1.05) {$x_{n+1}$};
\node[chico]    (ev2) at (0,-1.80) {$\vdots$};
\node[entrada]  (eN)  at (0,-2.55) {$x_N$};

% \node[notac,rotate=90] at (-1.35, 2.02) {primeras $n$\\ entradas};
% \node[notac,rotate=90] at (-1.35,-1.80) {entradas\\ restantes};

%=================== RED DE BAJA FIDELIDAD (BF) ========================
\foreach \j/\y in {1/3.07,2/2.02,3/0.97} {
  \node[neurona] (bf1-\j) at (2.6,\y) {};
  \node[neurona] (bf2-\j) at (4.1,\y) {};
}
\node[salida] (ubf) at (5.6,2.02) {$u_\text{LF}$};

% solo el subconjunto alimenta la red BF
\foreach \e in {e1,e2,en} {
  \foreach \j in {1,2,3} { \draw[syn] (\e) -- (bf1-\j); }
}
\foreach \i in {1,2,3} {
  \foreach \j in {1,2,3} { \draw[syn] (bf1-\i) -- (bf2-\j); }
  \draw[syn] (bf2-\i) -- (ubf);
}

\draw[grupoBF] (1.83,0.20) rectangle (6.37,3.84);
\node[etiqueta,anchor=south west] at (1.83,3.84)
  {Red de baja fidelidad $\theta_\text{LF}$ pre-entrenada};

%=========== BUS CON TODAS LAS ENTRADAS HACIA LA RED AF ================
% cada entrada, del subconjunto y de las restantes, se reparte al bus
\coordinate (nodoBus) at (2.6,-3.0);
\foreach \e in {e1,e2,en,er1,eN} { \draw[reparto] (\e) -- (nodoBus); }
\fill[black!75] (nodoBus) circle (1.6pt);

\draw[bus,rounded corners=3pt] (nodoBus) -- (5.9,-3.0) -- (5.9,-1.2) -- (7.30,-1.2);
\node[nota,anchor=north west] at (2.9,-3.4) {todas las entradas\\ $x_1,\dots,x_N$};

%=================== RED DE ALTA FIDELIDAD (AF) ========================
\node[barra] (concat) at (7.6,0) {};
\node[nota,rotate=90] at (7.6,0) {concatenación};

\foreach \j/\y in {1/1.70,2/0.57,3/-0.57,4/-1.70} {
  \node[neurona] (af1-\j) at (8.9,\y) {};
  \node[neurona] (af2-\j) at (10.4,\y) {};
}
\node[salida] (uaf) at (11.9,0) {$u_\text{HF}$};

\draw[flujo] (ubf.east) -- (7.30,1.3);
\foreach \j in {1,2,3,4} { \draw[syn] (concat.east) -- (af1-\j); }
\foreach \i in {1,2,3,4} {
  \foreach \j in {1,2,3,4} { \draw[syn] (af1-\i) -- (af2-\j); }
  \draw[syn] (af2-\i) -- (uaf);
}

\draw[grupoAF] (6.98,-2.52) rectangle (12.67,2.52);
\node[etiqueta,anchor=south west] at (6.98,2.52)
  {Red de alta fidelidad $\theta_\text{HF}$\\ informada por la física (PINN)};

%=============== RAMAS DE PÉRDIDA (FÍSICA Y DATOS) =====================
\node[bloque]  (ad)   at (4.6,-6.3)
  {Derivadas\\ automáticas\\ de $u_\text{HF}$ respecto\\ de las entradas};
\node[perdida] (lfis) at (4.6,-8.4)
  {$\mathcal{L}_{\text{fís}}=\big\|\mathcal{R}(u_\text{HF})\big\|^{2}$\\ residuo de la EDP};
\node[perdida,text width=3.4cm] (ldat) at (8.8,-6.3)
  {$\mathcal{L}_{\text{datos}}=\big\|u_\text{HF}-u^{*}\big\|^{2}$\\ $u^{*}$: datos medidos};
\node[perdida,text width=3.4cm] (ltot) at (8.8,-8.4)
  {$\mathcal{L}=\mathcal{L}_{\text{datos}}+\lambda\,\mathcal{L}_{\text{fís}}$};

% tronco desde u_AF, con derivacion hacia la rama de datos
\coordinate (bif) at (11.9,-4.9);
\draw[tronco] (uaf.south) -- (bif);
\draw[flujo,rounded corners=3pt] (bif) -- (4.6,-4.9) -- (ad.north);
\fill[black!75] (8.8,-4.9) circle (1.6pt);
\draw[flujo] (8.8,-4.9) -- (ldat.north);

\draw[flujo] (ad.south)   -- (lfis.north);
\draw[flujo] (ldat.south) -- (ltot.north);
\draw[flujo] (lfis.east)  -- (ltot.west);

%=================== RETROPROPAGACIÓN (SOLO RED AF) ====================
% vuelve por la derecha y entra al costado de la red AF, sin cruzar el tronco
\draw[bp,rounded corners=3pt]
  (ltot.east) -- (13.4,-8.4) -- (13.4,-1.8) -- (12.67,-1.8);
% \node[notac,anchor=north,text width=4.2cm] at (8.8,-9.3)
%   {retropropagación: sólo se\\ actualizan los pesos $\theta_{AF}$};

%============================ ACLARACIÓN ===============================
% \node[nota,anchor=north west,text width=4.0cm] at (-1.75,-4.9)
%   {la red BF sólo usa $x_1,\dots,x_n$; la red AF usa todas las entradas y
%    además $u_{BF}$};

\end{tikzpicture}
  }
  \caption{Esquema de la arquitectura de red neuronal de multifidelidad informada por física.}
  \label{fig:mfpinn}
\end{figure}
